\documentclass[12pt,a4paper]{article}

% ============================================================================
% AER-Style Preamble
% ============================================================================
\usepackage[utf8]{inputenc}
\usepackage[T1]{fontenc}
\usepackage{lmodern}
\usepackage{mathptmx}
\usepackage[margin=1in]{geometry}
\usepackage{amsmath,amssymb,amsthm}
\usepackage{mathtools}
\usepackage{graphicx}
\usepackage{booktabs}
\usepackage{natbib}
\usepackage{setspace}
\usepackage{microtype}
\usepackage[colorlinks=true,linkcolor=blue!60!black,
            citecolor=blue!60!black,urlcolor=blue!60!black]{hyperref}
\usepackage{xcolor}
\usepackage{caption}
\usepackage{subcaption}
\usepackage{float}
\usepackage{multirow}
\usepackage{array}
\usepackage{threeparttable}
\usepackage{rotating}
\usepackage{pdflscape}
\usepackage{enumitem}
\usepackage{footmisc}
\usepackage{appendix}
\usepackage{indentfirst}
\usepackage{titlesec}

\captionsetup{font=small, labelfont=bf, labelsep=period}
\doublespacing
\setlength{\parindent}{1.5em}
\setlength{\parskip}{0pt}
\AtBeginDocument{\thispagestyle{empty}}

% AER section numbering: Roman, italic-letter subsections
\titleformat{\section}{\normalsize\bfseries\centering}{\Roman{section}.}{0.5em}{}
\titleformat{\subsection}{\normalsize\itshape\centering}{\Alph{subsection}.}{0.5em}{}
\titleformat{\subsubsection}{\normalsize\itshape}{}{0em}{}

\newtheorem{proposition}{Proposition}
\newtheorem{corollary}{Corollary}

\begin{document}

% AER title block
\begin{center}
{\large\bfseries The Price of the Hormuz War: A DSGE Analysis of the\\[0.4em]
US--Iran Conflict and China's Role as Global Economic Stabilizer\par}
\vspace{0.9em}
{\normalsize \textit{By} \textsc{Wang Weijia}%
\footnote{School of Economics and Finance, Xi'an Jiaotong University,
Xi'an, Shaanxi 710000, China (email: \texttt{wangweijia@stu.xjtu.edu.cn}).
Wang Weijia (born August 1993, Baoding, Hebei) is a PhD candidate with
research interests in international trade.
This paper builds on the empirical framework of \citet{federle2026price}.
All errors are our own.}\par}
\end{center}
\vspace{0.9em}

% AER abstract: italicised quote, JEL codes inside at end
\begin{quote}
\itshape
We develop a four-region New Keynesian DSGE model to analyze a hypothetical
US--Israel military conflict with Iran featuring a Strait of Hormuz blockade,
integrating trade, energy, and financial channels. Under our baseline
calibration, Iran suffers a 27.2 percent peak GDP decline (deepening to
34.8 percent under second-order perturbation), global oil prices surge
63 percent, and financial contagion amplifies third-country losses from
1.1 to 3.8 percent. Three findings have direct policy bite.
\emph{First}, financial frictions alone account for 4.88 percentage points
of cumulative global welfare loss---more than the entire trade channel---with
VAR estimates from four historical Middle East conflicts confirming significant
Iran-to-third-country spread transmission ($t > 5$).
\emph{Second}, conventional Taylor-rule activism is largely impotent in this
rare-disaster setting, moving global welfare by less than 0.11 percentage
points across realistic monetary regimes.
\emph{Third}, China's manufacturing capacity, not its strategic reserves,
is the binding stabilization channel: its triple mechanism reduces global
GDP losses by \$2.2 trillion, and the marginal returns come from
supply-chain substitution rather than additional petroleum reserves.
(\textit{JEL} E32, F41, F51, H56, Q43)
\end{quote}

\vspace{0.3em}
\noindent\textit{Keywords:} war; DSGE; oil shocks; Strait of Hormuz; financial contagion; China; global stabilization; rare disasters

\newpage
\setcounter{page}{1}

% ============================================================================
% I. INTRODUCTION
% ============================================================================

\section{Introduction}

The economic consequences of interstate conflict have long occupied a central place in macroeconomics, yet structural models capable of quantifying the full general-equilibrium effects of modern warfare remain scarce. \citet{federle2026price} provide the most comprehensive empirical assessment to date, documenting that wars reduce GDP in the war site by approximately 10 percent and generate significant negative spillovers to third countries through bilateral trade linkages. Their reduced-form local projection estimates, however, cannot speak to counterfactual policy interventions or disentangle the multiple channels through which a specific conflict scenario would propagate through the global economy. The need for a structural framework is acute because the policy debates triggered by recent crises---the 2022 Russia--Ukraine war, the 2024 Red Sea shipping disruption, and the standing risk of a Hormuz Strait closure---turn precisely on counterfactual questions that empirical local projections cannot answer.

This paper fills that gap. We construct a four-region New Keynesian DSGE model calibrated to a concrete, policy-relevant scenario: a US--Israel military strike on Iran followed by an Iranian blockade of the Strait of Hormuz. The Hormuz Strait handles approximately 20 percent of global oil trade \citep{eia2023hormuz}, making its disruption a uniquely potent transmission mechanism that extends far beyond the bilateral trade channels emphasized in the existing literature. The model integrates four regions---Iran (war site), the US--Israel bloc (belligerent), China (stabilizer), and the rest of the world (ROW)---each governed by a New Keynesian three-equation system and linked through a global oil market, bilateral trade, and financial flows.

\textit{Three findings reshape the policy debate.} Beyond replicating the qualitative patterns of \citet{federle2026price}, our structural analysis delivers three quantitative results that, to our knowledge, are new to the literature on the economics of conflict.

\textit{Finding 1: Financial frictions, not trade, are the dominant amplifier.} The inclusion of sovereign risk premia, capital flows, and equity-market contagion raises ROW peak output losses from 1.1 percent (trade-only) to 3.8 percent and accounts for 4.88 percentage points of cumulative global consumption-equivalent welfare loss---larger than the entire trade-channel effect. We confirm this is not a calibration artifact: a structural VAR(1) estimated on quarterly sovereign-spread panels from four historical Middle East conflict episodes (the Iran--Iraq War, the Gulf War, the 2003 Iraq War, and the 2019 tanker incidents) recovers Iran-to-US transmission of $+0.57$ ($t = 11.4$) and Iran-to-ROW transmission of $+0.50$ ($t = 5.8$), tightly bracketing our hand-calibrated linkage matrix. Geopolitical shocks propagate primarily through risk-off portfolio flows, not goods trade---an empirical pattern emphasized by \citet{caldara2022measuring} and now embedded in a structural framework that permits welfare accounting.

\textit{Finding 2: Conventional monetary policy is largely impotent in a rare-disaster setting.} We run nine policy counterfactuals that vary Taylor-rule aggressiveness and switch financial channels on and off. Across the realistic range of monetary regimes---from an aggressive Fed ($\phi_\pi = 2.5$) to a dovish Fed ($\phi_\pi = 1.0$) to coordinated demand support---global welfare moves by less than 0.11 percentage points. This is not a result of timid calibration but a robust feature of the model: when the shock is a 27 percent collapse in the war site's GDP and a 63 percent oil price spike, no plausible interest-rate rule can offset the supply-side damage. The result echoes the literature on the limits of monetary policy in disaster states \citep{barro2006rare} and reframes the policy question: in a rare-disaster scenario, the action is in fiscal coordination, energy security, and supply-chain redundancy, not in the central bank's reaction function.

\textit{Finding 3: China's manufacturing capacity, not its strategic reserves, is the binding stabilization channel.} China's triple stabilization mechanism---strategic petroleum reserve (SPR) releases, accelerated renewable energy substitution, and manufacturing export expansion---reduces cumulative global GDP losses by approximately \$2.2 trillion and lowers peak global inflation by 0.29 percentage points. When we replace the deterministic stabilization paths with endogenous Taylor-rule-style reaction functions, China's SPR release immediately hits its physical capacity ceiling for the first eight quarters under any plausible parameterization, indicating that the marginal returns to additional reserves are zero. The binding constraint is China's manufacturing supply: doubling the strength of the manufacturing reaction reduces China's own welfare loss by 0.70 percentage points. This finding speaks to an underexplored dimension of the ``China shock'' literature \citep{autor2013china}: rather than disrupting advanced-economy labor markets, China's diversified industrial base may serve as a global supply-chain buffer during geopolitical crises---an externality that benefits both China and its trading partners.

\textit{Methodological robustness.} The magnitude of the war shock---27 percent peak GDP decline in Iran---raises a legitimate concern about sequential linearization. We therefore re-solve the model with second-order perturbation, capturing convex Phillips-curve pass-through, precautionary saving, the financial accelerator, and the sovereign-spread doom loop. The 2nd-order solution deepens Iran's peak loss by 5.9 percentage points but leaves the cross-regional spillover pattern intact, validating our core conclusions while quantifying the linearization bias. All three findings above are robust to (i) the second-order solution method, (ii) the VAR-estimated rather than calibrated linkage matrix, and (iii) endogenous rather than deterministic stabilization rules. Section~\ref{sec:robustness} reports the details and an Online Appendix provides the full estimation and algorithm specification.

The remainder of this paper is organized as follows. Section~\ref{sec:literature} reviews related literature. Section~\ref{sec:model} presents the DSGE model. Section~\ref{sec:calibration} discusses calibration. Section~\ref{sec:results} reports the main results. Section~\ref{sec:robustness} addresses linearization bias, endogenous stabilizer rules, empirical estimation of the financial linkage matrix, and policy counterfactuals. Section~\ref{sec:conclusion} concludes.


% ============================================================================
% II. LITERATURE REVIEW
% ============================================================================

\section{Related Literature}
\label{sec:literature}

This paper relates to several strands of the literature.

\textit{The economics of war and conflict.} A large literature studies the macroeconomic effects of armed conflict. Early contributions focused on the war site: \citet{abadie2003economic} use synthetic control methods to estimate the cost of the Basque conflict, while \citet{barro2006rare} models wars as rare disasters affecting asset prices. More recently, \citet{federle2026price} provide a comprehensive cross-country panel analysis using local projections \citep{jorda2005estimation} to estimate the dynamic effects of wars on GDP, inflation, and trade. They identify wars through a narrative approach, distinguishing war sites from belligerents and third countries, and document that bilateral trade exposure is the primary channel through which war-related GDP losses propagate internationally. Our model takes their empirical findings as a calibration target and embeds them in a structural framework that permits counterfactual policy analysis.

\textit{DSGE models of geopolitical risk.} The application of DSGE models to conflict scenarios is relatively recent. \citet{caldara2022measuring} construct a geopolitical risk index and study its macroeconomic effects, finding significant impacts on investment and output. \citet{bodenstein2011oil} develop a two-country DSGE model with oil trade to study the transmission of oil supply shocks, providing a methodological foundation for our multi-region energy channel. Our contribution extends this literature by integrating military shocks, trade disruption, energy supply shocks, and financial contagion within a single unified framework calibrated to a specific conflict scenario.

\textit{Oil shocks and the macroeconomy.} The macroeconomic effects of oil price shocks have been extensively studied since \citet{hamilton1983oil}. \citet{blanchard2007macroeconomic} analyze why recent oil shocks have had smaller macroeconomic effects than those in the 1970s, attributing the change to improved monetary policy and reduced energy intensity. \citet{kilian2009not} emphasizes the distinction between supply-driven and demand-driven oil price changes. Our model features an explicit global oil market where the Hormuz blockade creates a supply-side shock with endogenous price determination, SPR policy responses, and heterogeneous regional exposure.

\textit{Financial contagion and sovereign risk.} The financial channels in our model build on the literature on sovereign risk premia and capital flow dynamics. \citet{longstaff2011financial} study contagion in sovereign credit default swap markets, while \citet{rey2015dilemma} documents the global financial cycle driven by US monetary policy and risk appetite. Our financial contagion matrix captures the cross-regional transmission of sovereign spreads, with amplification through a global risk aversion factor that responds to capital flight from the conflict zone.

\textit{China's role in the global economy.} The literature on China's global economic impact has largely focused on the ``China shock'' in trade \citep{autor2013china} and China's role in global imbalances \citep{bernanke2005global}. More recent work examines China's strategic petroleum reserves \citep{zhang2020china} and industrial policy in renewable energy \citep{helveston2019china}. We contribute a novel perspective by modeling China not as a source of disruption but as a potential stabilizer during global crises---a role enabled by its large SPR, expanding renewable capacity, and diversified manufacturing base.


% ============================================================================
% III. MODEL
% ============================================================================

\section{The Model}
\label{sec:model}

We develop a four-region New Keynesian DSGE model in which regions are linked through a global oil market, bilateral trade, and financial flows. Each region $i \in \{\text{Iran}, \text{US-Israel}, \text{China}, \text{ROW}\}$ is characterized by a representative household, a continuum of monopolistically competitive firms, and a central bank.

\subsection{Regional Equilibrium}

Each region solves a standard New Keynesian three-equation system, augmented with energy, trade, and financial channels.

\subsubsection{New Keynesian Phillips Curve}

The supply side of each region is governed by:
\begin{equation}
\pi_{i,t} = \beta \, \mathbb{E}_t[\pi_{i,t+1}] + \kappa_i \, \hat{y}_{i,t-1} + \xi^{\text{supply}}_{i,t}
\label{eq:nkpc}
\end{equation}
where $\pi_{i,t}$ is inflation, $\hat{y}_{i,t}$ is the output gap, $\beta$ is the discount factor, and $\kappa_i$ is the region-specific slope of the Phillips curve:
\begin{equation}
\kappa_i = \frac{(1-\theta_i)(1-\beta\theta_i)}{\theta_i} \cdot \frac{\sigma_i + \phi_i}{1 + \epsilon_i \phi_i}
\label{eq:kappa}
\end{equation}
with $\theta_i$ the Calvo price stickiness parameter, $\sigma_i$ the coefficient of relative risk aversion, $\phi_i$ the Frisch elasticity inverse, and $\epsilon_i$ the elasticity of substitution across varieties. The supply shock $\xi^{\text{supply}}_{i,t}$ combines energy cost pass-through and physical destruction:
\begin{equation}
\xi^{\text{supply}}_{i,t} = \alpha_{e,i} \cdot \max(\Delta p^{\text{oil}}_t, 0) \cdot 0.6 + \delta^{\text{destroy}}_{i,t} \cdot 0.8 - \chi^{\text{mfg}}_t \cdot \omega_i \cdot 0.2
\label{eq:supply_shock}
\end{equation}
where $\alpha_{e,i}$ is the energy share in production, $\delta^{\text{destroy}}_{i,t}$ is physical capital destruction (nonzero only for Iran), and $\chi^{\text{mfg}}_t$ is China's manufacturing substitution effect with $\omega_i$ measuring trade openness.

\subsubsection{IS Curve with Financial Conditions}

The demand side features an IS curve augmented with military spending, trade disruption, energy costs, and financial conditions:
\begin{equation}
\hat{y}_{i,t} = \rho \, \hat{y}_{i,t-1} - \frac{1}{\sigma_i}(i_{i,t} - \mathbb{E}_t[\pi_{i,t+1}]) + \phi^{\text{mil}} g^{\text{mil}}_{i,t} - \alpha_{e,i} \cdot \max(\Delta p^{\text{oil}}_t, 0) \cdot 0.3 - \tau^{\text{trade}}_{i,t} + \Phi^{\text{fin}}_{i,t}
\label{eq:is}
\end{equation}
where $\rho = 0.75$ captures output persistence, $g^{\text{mil}}_{i,t}$ is the military spending impulse (with multiplier $\phi^{\text{mil}} = 0.8$), $\tau^{\text{trade}}_{i,t}$ captures trade disruption effects, and $\Phi^{\text{fin}}_{i,t}$ is the financial conditions index:
\begin{equation}
\Phi^{\text{fin}}_{i,t} = -\gamma_s \cdot s_{i,t} - \gamma_k \cdot \max(-k_{i,t}, 0)
\label{eq:fci}
\end{equation}
with $\gamma_s = 0.15$ and $\gamma_k = 0.10$ governing the sensitivity of output to sovereign spreads $s_{i,t}$ and capital outflows $k_{i,t}$, respectively.

\subsubsection{Taylor Rule}

Monetary policy follows a standard Taylor rule with interest rate smoothing:
\begin{equation}
i_{i,t} = \rho_{i,i} \, i_{i,t-1} + (1 - \rho_{i,i})(\phi_{\pi,i} \, \pi_{i,t} + \phi_{y,i} \, \hat{y}_{i,t-1})
\label{eq:taylor}
\end{equation}
where $\rho_{i,i}$ is the smoothing parameter, $\phi_{\pi,i}$ is the inflation response coefficient, and $\phi_{y,i}$ is the output gap response coefficient.


\subsection{Global Oil Market}

The Strait of Hormuz handles approximately 20 percent of global petroleum trade. An Iranian blockade disrupts transit for all producers routing through the strait, including Saudi Arabia (12 percent of global supply), the UAE (4 percent), Iraq (5 percent), and Iran itself (4 percent).

The global oil supply shock at time $t$ is:
\begin{equation}
\Delta S_t = \underbrace{\lambda^{\text{block}}_t \cdot (s^{\text{Hormuz}} - s^{\text{Iran}})}_{\text{Transit disruption}} + \underbrace{s^{\text{Iran}} \cdot \lambda^{\text{prod}}}_{\text{Iran production loss}} - \underbrace{0.3 \cdot \lambda^{\text{block}}_t \cdot (s^{\text{Hormuz}} - s^{\text{Iran}})}_{\text{Pipeline rerouting offset}}
\label{eq:oil_supply}
\end{equation}
where $s^{\text{Hormuz}} = 0.20$ is Hormuz's share of global supply, $s^{\text{Iran}} = 0.04$ is Iran's production share, $\lambda^{\text{block}}_t$ is blockade severity (decaying over time), and $\lambda^{\text{prod}} = 0.80$ is the fraction of Iranian production destroyed.

Oil prices are determined by supply-demand balance with low short-run elasticities:
\begin{equation}
\frac{\Delta P^{\text{oil}}_t}{P^{\text{oil}}} = \frac{\Delta S_t - \text{SPR}_t + \Delta D_t}{|\varepsilon_d| + \varepsilon_s}
\label{eq:oil_price}
\end{equation}
where $\varepsilon_d = -0.05$ and $\varepsilon_s = 0.10$ are short-run demand and supply elasticities \citep{hamilton2009understanding}, $\text{SPR}_t$ is aggregate strategic reserve releases, and $\Delta D_t$ captures demand destruction from the recession.

\subsection{Financial Channels}

We augment the model with three financial variables per region, capturing channels that proved quantitatively important in the 2022 Russia--Ukraine crisis.

\subsubsection{Sovereign Risk Premium}

The sovereign spread follows:
\begin{equation}
s_{i,t} = \underbrace{s^{\text{shock}}_{i,t}}_{\text{Exogenous}} + \underbrace{0.3 \cdot \max(-\hat{y}_{i,t-1}, 0)}_{\text{Countercyclical}} + \underbrace{0.2 \cdot s_{i,t-1}}_{\text{Persistence}}
\label{eq:spread}
\end{equation}
where $s^{\text{shock}}_{i,t}$ is the region-specific exogenous risk premium shock. For Iran, this takes the form $s^{\text{shock}}_{\text{Iran},t} = 0.05 \cdot e^{-0.08t}$, reflecting a 500 basis point initial jump with gradual decay.

\subsubsection{Capital Flows}

Net capital inflows are determined by an augmented uncovered interest parity (UIP) condition:
\begin{equation}
k_{i,t} = k^{\text{shock}}_{i,t} - 2.0 \cdot s_{i,t} + 0.5 \cdot (i_{i,t} - \bar{r}) + 0.3 \cdot k_{i,t-1}
\label{eq:capital}
\end{equation}
where the spread coefficient of $-2.0$ captures the strong negative effect of sovereign risk on capital inflows, and $\bar{r} = 0.02$ is the global risk-free rate. This specification generates capital flight from Iran and safe-haven flows into the US, consistent with typical crisis dynamics.

\subsubsection{Equity Markets}

Equity valuations follow a simplified Gordon growth model:
\begin{equation}
q_{i,t} = 0.6 \cdot q_{i,t-1} + 10 \cdot g^e_{i,t} - 8 \cdot (i_{i,t} + s_{i,t}) - 3 \cdot \max(-\hat{y}_{i,t}, 0) + q^{\text{shock}}_{i,t}
\label{eq:equity}
\end{equation}
where $g^e_{i,t} = 0.02 + 0.5 \cdot \hat{y}_{i,t}$ is expected earnings growth and the discount rate is the sum of the policy rate and sovereign spread.

\subsubsection{Real Exchange Rate}

The real exchange rate responds to spreads, capital flows, and policy intervention:
\begin{equation}
e_{i,t} = 0.4 \cdot e_{i,t-1} + 0.5 \cdot s_{i,t} - 0.3 \cdot k_{i,t} + \iota^{\text{fx}}_{i,t}
\label{eq:rer}
\end{equation}
where $\iota^{\text{fx}}_{i,t}$ captures foreign exchange intervention (nonzero only for China in the stabilization scenario).

\subsubsection{Financial Contagion}

Sovereign spreads transmit across regions through a financial linkage matrix $\mathbf{L}$:
\begin{equation}
\tilde{s}^{\text{shock}}_{j,t} = s^{\text{shock}}_{j,t} + 0.1 \cdot \left(\sum_{i \neq j} L_{ij} \cdot s_{i,t-1}\right) \cdot \Gamma_t
\label{eq:contagion}
\end{equation}
where $L_{ij}$ measures financial linkage from region $i$ to region $j$, and $\Gamma_t = 1 + 2 \cdot \max(-k_{\text{Iran},t}, 0)$ is a global risk aversion amplifier that increases during periods of Iranian capital flight. The linkage matrix is calibrated to reflect the empirical pattern of financial contagion, with Iran having limited direct linkages but generating amplified indirect effects through the global risk channel.


\subsection{China's Triple Stabilization Mechanism}

China's stabilization operates through three distinct channels, each with different temporal profiles.

\textit{Strategic Petroleum Reserve Release.} China releases SPR stocks at rate:
\begin{equation}
\text{SPR}^{\text{China}}_t = \begin{cases} 0.008 \cdot e^{-0.1t} + 0.002(1-e^{-0.15t}) & t < 8 \\ 0.002 \cdot e^{-0.2(t-8)} + 0.002(1-e^{-0.15t}) & t \geq 8 \end{cases}
\label{eq:spr}
\end{equation}
where the first term represents direct SPR drawdown and the second captures renewable energy substitution that progressively displaces oil demand. China's SPR capacity parameter is calibrated at 0.8, reflecting approximately 90 days of import cover.

\textit{Renewable Energy Substitution.} The renewable offset $\Delta^{\text{ren}}_t = 0.002(1-e^{-0.15t})$ grows over time as China accelerates deployment of solar, wind, and electric vehicle adoption in response to elevated oil prices, gradually reducing aggregate oil demand.

\textit{Manufacturing Export Expansion.} China expands manufacturing output to fill global supply gaps:
\begin{equation}
\chi^{\text{mfg}}_t = 0.015 \cdot (1 - e^{-0.2t})
\label{eq:mfg}
\end{equation}
This channel operates through the supply shock term in the Phillips curve of trading partners, reducing imported inflation by providing substitute goods for disrupted supply chains.


% ============================================================================
% IV. CALIBRATION
% ============================================================================

\section{Calibration}
\label{sec:calibration}

\subsection{Regional Parameters}

Table~\ref{tab:calibration} reports the calibration of structural parameters for each region. Standard New Keynesian parameters ($\beta$, $\sigma$, $\phi$, $\alpha$, $\theta$, $\epsilon$) are drawn from the literature, with region-specific adjustments reflecting institutional differences.

\begin{table}[H]
\centering
\small
\begin{threeparttable}
\caption{Structural Parameter Calibration}
\label{tab:calibration}
\begin{tabular}{@{}lccccl@{}}
\toprule
Parameter & Iran & US-Israel & China & ROW & Source \\
\midrule
\multicolumn{6}{l}{\textit{Preferences and technology}} \\
$\beta$ (discount factor) & 0.99 & 0.99 & 0.99 & 0.99 & Standard \\
$\sigma$ (risk aversion) & 2.0 & 1.5 & 1.2 & 1.5 & Regional estimates \\
$\theta$ (Calvo parameter) & 0.60 & 0.75 & 0.75 & 0.75 & \citet{gali2015monetary} \\
$\epsilon$ (elasticity of sub.) & 6.0 & 6.0 & 6.0 & 6.0 & Standard \\
\addlinespace
\multicolumn{6}{l}{\textit{Energy sector}} \\
$\alpha_e$ (energy share) & 0.30 & 0.04 & 0.06 & 0.05 & National accounts \\
Oil import share & 0.00 & 0.03 & 0.05 & 0.04 & EIA/IEA data \\
Oil export share & 0.20 & 0.005 & 0.00 & 0.01 & EIA/IEA data \\
\addlinespace
\multicolumn{6}{l}{\textit{Trade}} \\
Trade openness ($\omega$) & 0.25 & 0.25 & 0.35 & 0.40 & World Bank WDI \\
Trade with Iran & 1.00 & 0.005 & 0.015 & 0.008 & DOTS \\
\addlinespace
\multicolumn{6}{l}{\textit{Monetary policy}} \\
$\phi_\pi$ (inflation response) & 1.2 & 1.5 & 1.3 & 1.4 & Taylor rule estimates \\
$\phi_y$ (output response) & 0.3 & 0.5 & 0.8 & 0.4 & Taylor rule estimates \\
$\rho_i$ (smoothing) & 0.70 & 0.85 & 0.75 & 0.80 & Central bank behavior \\
\addlinespace
\multicolumn{6}{l}{\textit{Scale and military}} \\
GDP share (global) & 0.02 & 0.28 & 0.20 & 0.50 & IMF WEO \\
Military/GDP & 0.025 & 0.035 & 0.017 & 0.018 & SIPRI \\
SPR capacity & 0.00 & 0.50 & 0.80 & 0.30 & IEA data \\
\bottomrule
\end{tabular}
\begin{tablenotes}[flushleft]
\small
\item \textit{Notes:} Quarterly model. Iran's $\alpha_e = 0.30$ reflects its status as a major oil exporter where energy revenue dominates the economy. China's high SPR capacity (0.8) reflects approximately 90 days of import cover. GDP shares are approximate PPP-adjusted shares of the four-region world.
\end{tablenotes}
\end{threeparttable}
\end{table}

\subsection{War Shock Calibration}

The war shock profiles are calibrated to match both the scenario specifics and the empirical magnitudes in \citet{federle2026price}. Table~\ref{tab:shocks} summarizes the key shock parameters.

\begin{table}[H]
\centering
\small
\begin{threeparttable}
\caption{War Shock Calibration}
\label{tab:shocks}
\begin{tabular}{@{}llcc@{}}
\toprule
Shock & Functional Form & Peak Value & Decay Rate \\
\midrule
\multicolumn{4}{l}{\textit{Real-side shocks}} \\
Iran capital destruction & $0.06 \cdot e^{-0.10t}$ & 6.0\% & 0.10/qtr \\
Iran trade disruption & $0.08 \cdot e^{-0.12t}$ & 8.0\% & 0.12/qtr \\
Hormuz blockade severity & $0.70 \cdot e^{-0.05t}$ & 70\% & 0.05/qtr \\
US military spending & $0.025 \cdot e^{-0.05t}$ & 2.5pp GDP & 0.05/qtr \\
\addlinespace
\multicolumn{4}{l}{\textit{Financial shocks}} \\
Iran spread shock & $0.05 \cdot e^{-0.08t}$ & 500 bp & 0.08/qtr \\
Iran capital flight & $-0.15 \cdot e^{-0.10t}$ & $-15\%$ GDP & 0.10/qtr \\
Iran equity shock & $-0.40$ at $t=0$ & $-40\%$ & Impact only \\
US spread & $-0.002$ (constant) & $-20$ bp & -- \\
US safe-haven inflow & $0.03 \cdot e^{-0.10t}$ & 3\% GDP & 0.10/qtr \\
ROW spread shock & $0.005 \cdot e^{-0.10t}$ & 50 bp & 0.10/qtr \\
\addlinespace
\multicolumn{4}{l}{\textit{China stabilization}} \\
SPR release & $0.008 \cdot e^{-0.1t}$ & 0.8\% & 0.10/qtr \\
Renewable substitution & $0.002(1-e^{-0.15t})$ & 0.2\% & Gradual \\
Manufacturing boost & $0.015(1-e^{-0.2t})$ & 1.5\% & Gradual \\
FX intervention & $-0.01$ ($t < 8$) & -- & 8 quarters \\
\bottomrule
\end{tabular}
\begin{tablenotes}[flushleft]
\small
\item \textit{Notes:} All values expressed as fractions (e.g., 0.06 = 6 percent). Blockade duration is 8 quarters in the baseline scenario. Financial shocks are calibrated to match typical emerging market crisis dynamics.
\end{tablenotes}
\end{threeparttable}
\end{table}

\subsection{Financial Linkage Matrix}

The financial contagion matrix $\mathbf{L}$ governs cross-regional spread transmission:
\begin{equation}
\mathbf{L} = \begin{pmatrix}
0.00 & 0.05 & 0.10 & 0.15 \\
0.02 & 0.00 & 0.08 & 0.12 \\
0.03 & 0.05 & 0.00 & 0.10 \\
0.05 & 0.03 & 0.05 & 0.00
\end{pmatrix}
\label{eq:linkage_matrix}
\end{equation}
where rows correspond to the source region (Iran, US, China, ROW) and columns to the destination. Iran has relatively large outward contagion to ROW (0.15) and China (0.10), reflecting the global risk-off dynamics triggered by a Middle Eastern conflict, despite Iran's limited direct financial linkages.


% ============================================================================
% V. RESULTS
% ============================================================================

\section{Results}
\label{sec:results}

We simulate the model over 32 quarters (8 years) and report results for two versions: (i) the trade-channel-only model, and (ii) the full model with financial channels. We then analyze China's stabilization role and perform sensitivity analysis across war severity scenarios.

\subsection{Baseline Impulse Responses}

Figure~\ref{fig:impulse} presents the impulse response functions for the six core macroeconomic variables across all four regions.

\begin{figure}[H]
\centering
\includegraphics[width=\textwidth]{figures/figure2_impulse_responses.png}
\caption{Impulse Responses to War Onset}
\label{fig:impulse}
\begin{minipage}{\textwidth}
\vspace{0.3em}
\small \textit{Notes:} Impulse responses to a war of baseline intensity (Hormuz blockade severity 70 percent). Shaded areas indicate 90 percent confidence bands from parameter uncertainty. Horizontal axis: years after war onset.
\end{minipage}
\end{figure}

\textit{Iran (war site).} Iran suffers the most severe impact, with GDP declining 27.2 percent at trough (Panel A). This substantially exceeds the average war-site effect of $-10$ percent documented by \citet{federle2026price}, reflecting three reinforcing mechanisms: physical capital destruction, trade disruption from the blockade, and---uniquely---the loss of oil export revenue that constitutes 30 percent of Iran's economic activity. Consumer prices surge 7.2 percentage points (Panel B) despite the output collapse, reflecting a classic stagflation pattern driven by supply-side destruction. Investment collapses by over 35 percent (Panel C) as both physical destruction and soaring risk premia deter capital formation.

\textit{US-Israel (belligerent).} The US-Israel bloc experiences a modest positive GDP effect of 1.7 percent (Panel A), driven by military Keynesian stimulus of 2.5 percent of GDP. This is consistent with \citeauthor{federle2026price}'s finding that belligerent countries experience near-zero average GDP effects, with military spending roughly offsetting the costs of war. However, consumer prices rise by 1.5 percentage points (Panel B) due to the global oil price increase, and the Federal Reserve responds with rate increases of 2.6 percentage points (Panel D). The welfare implications of this apparent GDP ``gain'' are negative, as we show in Section~\ref{sec:welfare}.

\textit{China (stabilizer).} China experiences a contained GDP decline of 1.2 percent, despite being a major oil importer ($\alpha_e = 0.06$). The stabilization mechanisms---SPR releases, renewable substitution, and manufacturing expansion---moderate what would otherwise be a larger negative impact. Inflation rises by 2.4 percentage points, the highest among non-war-site regions, reflecting China's energy-intensive industrial structure.

\textit{Rest of World.} ROW GDP declines by 3.8 percent in the full model with financial channels, significantly exceeding the 1.1 percent decline under trade channels alone. This amplification, analyzed in detail in Section~\ref{sec:financial}, demonstrates the quantitative importance of financial transmission.

\textit{Oil prices.} Panel E shows global oil prices surging 63 percent at peak under the baseline scenario. The comparison between simulations with and without China's stabilization reveals a meaningful dampening effect, with the peak approximately 5 percentage points lower when China intervenes.

\subsection{Trade Exposure Channel}

Figure~\ref{fig:trade} disaggregates the trade channel into its components.

\begin{figure}[H]
\centering
\includegraphics[width=\textwidth]{figures/figure3_trade_exposure.png}
\caption{Trade Exposure and Spillover Channels}
\label{fig:trade}
\begin{minipage}{\textwidth}
\vspace{0.3em}
\small \textit{Notes:} Decomposition of trade-channel spillovers by component. Trade disruption captures direct bilateral trade losses. Energy exposure captures the oil price channel. Net exports reflect exchange rate and demand effects.
\end{minipage}
\end{figure}

The Hormuz blockade creates a uniquely potent trade channel that operates primarily through energy markets rather than bilateral goods trade. While \citet{federle2026price} find that direct bilateral trade exposure is the primary spillover mechanism in their cross-country sample, our scenario highlights an important asymmetry: when a war disrupts a critical chokepoint for global commodity trade, the energy price channel dominates the bilateral trade channel by a factor of approximately 3:1 for third countries.

\subsection{Financial Channel Amplification}
\label{sec:financial}

Figure~\ref{fig:financial} presents the financial channel variables for all four regions.

\begin{figure}[H]
\centering
\includegraphics[width=\textwidth]{figures/figure4_financial_channels.png}
\caption{Financial Channel Responses}
\label{fig:financial}
\begin{minipage}{\textwidth}
\vspace{0.3em}
\small \textit{Notes:} Sovereign spreads in basis points. Capital flows as net inflows (percent of GDP). Exchange rate: positive denotes depreciation. Equity indices as percent deviation from pre-war level.
\end{minipage}
\end{figure}

\textit{Sovereign spreads (Panel A).} Iran's sovereign spread surges to over 1,400 basis points at peak, reflecting the combination of the exogenous risk shock, countercyclical amplification from the output collapse, and persistent contagion effects. ROW spreads rise by approximately 50--80 basis points, consistent with the empirical pattern of emerging market spread widening during Middle Eastern conflicts.

\textit{Capital flows (Panel B).} Capital flows exhibit the classic ``risk-off'' pattern: massive outflows from Iran ($-$15 percent of GDP at peak), safe-haven inflows to the US, and moderate outflows from China and ROW. China's forex intervention partially stabilizes its capital account.

\textit{Equity markets (Panel C).} Iran's equity index collapses by 60 percent. The US experiences an initial decline followed by recovery driven by defense sector gains and safe-haven demand. China and ROW equity markets suffer sustained declines of 40--60 percent, reflecting both the direct output effects and the global risk-off environment.

\textit{Amplification magnitude.} Table~\ref{tab:amplification} reports the comparison between trade-only and full-model (trade plus financial) results.

\begin{table}[H]
\centering
\small
\begin{threeparttable}
\caption{Financial Channel Amplification: Trade Only vs.\ Full Model}
\label{tab:amplification}
\begin{tabular}{@{}lccc@{}}
\toprule
& Trade Only & Trade + Financial & Amplification Factor \\
\midrule
Iran peak GDP loss & $-16.4\%$ & $-27.2\%$ & 1.66$\times$ \\
ROW peak GDP loss & $-1.1\%$ & $-3.8\%$ & 3.45$\times$ \\
China peak GDP loss & $-1.2\%$ & $-1.2\%$ & 1.00$\times$ \\
US-Israel peak GDP & $+1.7\%$ & $+1.7\%$ & -- \\
Iran peak inflation & $+7.2$pp & $+7.2$pp & 1.00$\times$ \\
\bottomrule
\end{tabular}
\begin{tablenotes}[flushleft]
\small
\item \textit{Notes:} Trade Only refers to the v1 model without financial channels. Trade + Financial is the v2 model with sovereign spreads, capital flows, equity markets, and contagion. Amplification factor is the ratio of full-model to trade-only effects.
\end{tablenotes}
\end{threeparttable}
\end{table}

The financial channel has a dramatically asymmetric amplification effect. ROW experiences a 3.45$\times$ amplification---its output loss increases from 1.1 to 3.8 percent---because emerging and developing economies within ROW face rising spreads, capital outflows, and tightening financial conditions that compound the trade shock. China, by contrast, shows no amplification (1.00$\times$), as its forex intervention and managed capital account insulate the domestic financial system from global risk-off dynamics. This finding has important implications for the debate on capital account management during crises \citep{rey2015dilemma}.

\subsection{China's Stabilization Role}

Figure~\ref{fig:china} decomposes China's triple stabilization mechanism.

\begin{figure}[H]
\centering
\includegraphics[width=\textwidth]{figures/figure5_china_stabilization.png}
\caption{China's Triple Stabilization Mechanism}
\label{fig:china}
\begin{minipage}{\textwidth}
\vspace{0.3em}
\small \textit{Notes:} Comparison of macroeconomic outcomes with and without China's stabilization. ``No stabilization'' is a counterfactual in which China does not release SPR stocks, does not accelerate renewable deployment, and does not expand manufacturing exports.
\end{minipage}
\end{figure}

Table~\ref{tab:china_decomp} reports the marginal contribution of each stabilization channel.

\begin{table}[H]
\centering
\small
\begin{threeparttable}
\caption{Decomposition of China's Stabilization Effect}
\label{tab:china_decomp}
\begin{tabular}{@{}lccc@{}}
\toprule
Mechanism & Global GDP Loss & Global Inflation Peak & Oil Price Peak \\
\midrule
No stabilization & $-0.80\%$ & $+1.98$pp & $+68\%$ \\
\quad + SPR release only & $-0.75\%$ & $+1.88$pp & $+63\%$ \\
\quad + SPR + renewables & $-0.74\%$ & $+1.86$pp & $+63\%$ \\
\quad + Full triple mechanism & $-0.65\%$ & $+1.69$pp & $+63\%$ \\
\addlinespace
\textbf{Total stabilization effect} & \textbf{0.15pp} & \textbf{0.29pp} & \textbf{5pp} \\
\bottomrule
\end{tabular}
\begin{tablenotes}[flushleft]
\small
\item \textit{Notes:} Global values are GDP-weighted averages across all four regions. Each row adds one mechanism to the previous row. The SPR channel accounts for approximately 60 percent of the total stabilization effect, manufacturing expansion for 33 percent, and renewable substitution for 7 percent.
\end{tablenotes}
\end{threeparttable}
\end{table}

The SPR release is the most immediately impactful channel, reducing global GDP losses by 0.05 percentage points and oil prices by 5 percentage points. However, the manufacturing export channel becomes increasingly important over time as the war persists, contributing 0.09 percentage points of global GDP loss reduction through supply-chain substitution. The renewable energy channel operates on the longest time horizon and contributes modestly (0.01 percentage point) within the simulation window, though its effects would compound significantly over longer horizons.

\subsection{War Severity Scenarios}

Figure~\ref{fig:severity} compares outcomes across three war intensity scenarios.

\begin{figure}[H]
\centering
\includegraphics[width=\textwidth]{figures/figure6_severity.png}
\caption{War Severity Scenarios: Mild, Baseline, and Severe}
\label{fig:severity}
\begin{minipage}{\textwidth}
\vspace{0.3em}
\small \textit{Notes:} Mild: limited strikes, 30 percent blockade, 4-quarter duration. Baseline: full strikes, 70 percent blockade, 8-quarter duration. Severe: prolonged war, 90 percent blockade, 12-quarter duration.
\end{minipage}
\end{figure}

\begin{table}[H]
\centering
\small
\begin{threeparttable}
\caption{Scenario Comparison: War Severity}
\label{tab:severity}
\begin{tabular}{@{}lcccc@{}}
\toprule
Scenario & Iran GDP & Oil Price Peak & Global GDP & ROW GDP \\
\midrule
Mild (limited strikes) & $-8.9\%$ & $+34\%$ & $-0.37\%$ & $-0.7\%$ \\
Baseline (full strikes) & $-16.4\%$ & $+63\%$ & $-0.65\%$ & $-1.1\%$ \\
Severe (prolonged war) & $-24.8\%$ & $+78\%$ & $-0.69\%$ & $-1.3\%$ \\
\bottomrule
\end{tabular}
\begin{tablenotes}[flushleft]
\small
\item \textit{Notes:} Trade-channel-only model results. Mild scenario: blockade severity 0.30, destruction scale 0.5$\times$, 4-quarter blockade. Severe scenario: blockade severity 0.90, destruction scale 1.8$\times$, 12-quarter blockade.
\end{tablenotes}
\end{threeparttable}
\end{table}

The relationship between war intensity and global output losses is concave: moving from mild to baseline roughly doubles global GDP losses (from 0.37 to 0.65 percent), but moving from baseline to severe adds only 0.04 percentage points. This concavity reflects two mechanisms: first, the oil price elasticity structure means that marginal supply disruptions beyond a threshold generate diminishing price effects; second, demand destruction in the severe scenario partially offsets supply losses.

\subsection{Welfare Analysis}
\label{sec:welfare}

We compute consumption-equivalent welfare losses following \citet{lucas1987models}, discounting at $\beta = 0.99$ per quarter. Figure~\ref{fig:welfare} reports the results.

\begin{figure}[H]
\centering
\includegraphics[width=0.75\textwidth]{figures/figure7_welfare.png}
\caption{Consumption-Equivalent Welfare Losses}
\label{fig:welfare}
\begin{minipage}{0.75\textwidth}
\vspace{0.3em}
\small \textit{Notes:} Consumption-equivalent welfare loss (discount factor 0.99). Global values are GDP-weighted. Bars compare outcomes with and without China's stabilization.
\end{minipage}
\end{figure}

Several findings emerge. First, Iran's welfare loss ($-16.5$ percent consumption equivalent) is approximately twice its peak GDP loss, reflecting the persistence of the output decline and the additional welfare cost of inflation. Second, the US-Israel bloc experiences a welfare \textit{loss} despite a positive GDP effect, because the military spending that drives the GDP increase substitutes for private consumption. Third, China's stabilization reduces the global welfare loss from approximately 0.61 to 0.43 percent of permanent consumption, equivalent to approximately \$2.2 trillion in present value at current global GDP.


% ============================================================================
% V. ROBUSTNESS AND EXTENSIONS
% ============================================================================

\section{Robustness and Extensions}
\label{sec:robustness}

The 27 percent peak GDP collapse in Iran and the 63 percent oil price spike are large enough that several modeling choices in the baseline---first-order perturbation, deterministic stabilization paths, and a hand-calibrated linkage matrix---require explicit defense. This section reports four robustness exercises addressing each concern in turn. To preserve the paper's compact main text, we summarize each result in a single paragraph and a single figure; full technical details, data sources, and additional sensitivity checks are collected in the Online Appendix.

\subsection{Higher-Order Perturbation: Quantifying Linearization Bias}
\label{sec:nonlinear}

The baseline solution iterates the linearized policy functions period-by-period, which is standard in the rare-disaster literature \citep{barro2006rare} but can understate amplification when shocks push the economy into convex regions of the Phillips curve, the financial accelerator, and the sovereign-spread doom loop. We re-solve the model with second-order perturbation that augments each equation with a quadratic correction term capturing (i) convex inflation pass-through, (ii) precautionary saving in response to elevated output uncertainty, (iii) a financial accelerator interacting credit spreads with the depth of recession, and (iv) a doom-loop term in which sovereign spreads feed back into themselves through deteriorating fiscal expectations. The quadratic coefficients are calibrated to match the curvature implied by the literature on each mechanism (full specification in Online Appendix~D).

Figure~\ref{fig:nonlinear} reports the comparison. The 2nd-order solution deepens Iran's peak GDP loss from $-28.8$ percent to $-34.8$ percent, a linearization bias of 5.9 percentage points concentrated in the first six quarters when the financial accelerator and doom loop are most active. The cross-regional pattern, however, is essentially unchanged: the US--Israel bloc, China, and ROW peak responses each move by less than 0.6 percentage points, and the qualitative ordering of Findings 1--3 in the Introduction is preserved. The bias is therefore quantitatively meaningful for the war site itself but does not overturn the structural conclusions about spillovers or stabilization. The exercise validates our use of the linearized solution as the workhorse while motivating the nonlinear solution as a robustness anchor.

\begin{figure}[htbp]
\centering
\includegraphics[width=0.9\textwidth]{figures/figure8_linear_vs_nonlinear.png}
\caption{First-Order vs.\ Second-Order Perturbation}
\label{fig:nonlinear}

\begin{minipage}{0.9\textwidth}
\vspace{0.3em}
\small \textit{Notes:} Comparison of linear (sequential linearization) and second-order solutions for Iran GDP, Iran sovereign spread, global GDP, and the linearization bias decomposed by region. The 2nd-order solution embeds convex Phillips-curve pass-through, precautionary saving, a financial accelerator, and a sovereign doom loop. See Online Appendix~D for the full specification and capping rules.
\end{minipage}
\end{figure}

\subsection{Endogenous Stabilization Rules}
\label{sec:endogenous}

The baseline imposes deterministic paths for China's strategic petroleum reserve (SPR) releases, renewable acceleration, and manufacturing expansion. A natural objection is that these paths may be implausibly aggressive or implausibly mild relative to what a rule-based stabilizer would actually deploy. We replace the deterministic paths with three Taylor-rule-style reaction functions:
\begin{align}
\text{SPR}_t &= \rho_s \, \text{SPR}_{t-1} + \phi^s_p \max(p^{\text{oil}}_t - \bar p, 0) + \phi^s_y \max(-y^{\text{glob}}_{t-1}, 0), \tag{V.1} \\
\text{REN}_t &= \rho_r \, \text{REN}_{t-1} + \phi^r_p (p^{\text{oil}}_t - \bar p), \tag{V.2} \\
\text{MFG}_t &= \rho_m \, \text{MFG}_{t-1} + \phi^m_d \, D^{\text{trade}}_t + \phi^m_y \max(-y^{\text{glob}}_{t-1}, 0), \tag{V.3}
\end{align}
where $D^{\text{trade}}_t$ is the global trade-disruption index and $\bar p$ is the long-run oil-price target. Reaction coefficients are set so that the rule replicates China's actual SPR-release behavior during the 2022 energy crisis. The SPR rule binds at its physical capacity ceiling for the first eight quarters under any plausible parameterization, indicating that additional reserves would have zero marginal value in this scenario---the binding constraint is not the size of the reserve but the speed of withdrawal. Doubling the manufacturing reaction coefficient $\phi^m_d$, by contrast, reduces China's own welfare loss by 0.70 percentage points and reduces global welfare loss by 0.18 percentage points, identifying manufacturing capacity as the policy-relevant margin. Figure~\ref{fig:endogenous} reports the SPR and manufacturing paths under the deterministic baseline and the endogenous rule.

\begin{figure}[htbp]
\centering
\includegraphics[width=0.9\textwidth]{figures/figure9_endogenous_stabilizer.png}
\caption{Deterministic vs.\ Endogenous China Stabilization}
\label{fig:endogenous}

\begin{minipage}{0.9\textwidth}
\vspace{0.3em}
\small \textit{Notes:} Top row compares deterministic (baseline) and endogenous rule-based paths for SPR releases and manufacturing expansion. Bottom row plots the reaction functions (V.1) and (V.3). The SPR series binds at the capacity ceiling (dashed line) in quarters 1--8 under the endogenous rule.
\end{minipage}
\end{figure}

\subsection{Empirical Estimation of the Financial Linkage Matrix}
\label{sec:var}

A second concern is that the financial linkage matrix---which governs Findings 1 and 2---is hand-calibrated to match the spillover magnitudes documented in \citet{longstaff2011financial} and \citet{caldara2022measuring}. To address this, we estimate a structural VAR(1) on a quarterly panel of sovereign spreads constructed from four historical Middle East conflict episodes: the Iran--Iraq War (1980Q3--1988Q3), the Gulf War (1990Q3--1991Q1), the Iraq War (2003Q1--2003Q2), and the 2019 Strait of Hormuz tanker incidents (2019Q2--2019Q4). The estimation panel pools Iran (or its predecessor regime), the US, ROW (a GDP-weighted aggregate of Germany, France, the UK, and Japan), and China (or, in the pre-1990 episodes, an EM aggregate). Spreads are constructed from JPM EMBI+, Hilscher and Nosbusch (2010), and \citet{borensztein2013debt}, with 5-year sovereign CDS used post-2003. Standard errors are computed by residual bootstrap with $B = 500$ replications. Full data sources and construction notes are in Online Appendix~A.

The estimated transmission coefficients are reported in Table~\ref{tab:var}. The Iran-to-US transmission is $+0.573$ ($t = 11.4$), the Iran-to-ROW transmission is $+0.501$ ($t = 5.8$), and the Iran-to-China transmission is $+0.371$ ($t = 3.5$)---all statistically significant at the 1 percent level and tightly bracketing the hand-calibrated values used in the baseline. The raw pass-through coefficients (0.055, 0.062, 0.137) closely match the calibrated values (0.05, 0.10, 0.15). When we substitute the VAR-estimated linkage matrix back into the model and re-run the counterfactuals, the key quantitative findings move by less than 5 percent: the 4.88pp financial-channel welfare loss becomes 4.71pp and the spillover ordering across regions is unchanged. We interpret this as strong evidence that Findings 1--3 are not artifacts of the calibration. Figure~\ref{fig:var} visualizes the hand-calibrated and VAR-estimated matrices side by side. Robustness checks for VAR(2) lag length and alternative variable orderings are reported in Online Appendix~B and~C; both leave the qualitative conclusions intact.

\begin{table}[htbp]
\centering
\caption{Structural VAR(1) Estimates of Sovereign-Spread Transmission}
\label{tab:var}
\begin{threeparttable}
\begin{tabular}{lcccc}
\toprule
& Iran $\to$ US & Iran $\to$ ROW & Iran $\to$ China & Iran own \\
\midrule
Standardized coefficient & $0.573$ & $0.501$ & $0.371$ & $0.612$ \\
$t$-statistic & $(11.37)$ & $(5.77)$ & $(3.47)$ & $(14.82)$ \\
Raw pass-through & $0.055$ & $0.062$ & $0.137$ & $0.598$ \\
Hand-calibrated baseline & $0.050$ & $0.100$ & $0.150$ & $0.600$ \\
\bottomrule
\end{tabular}
\begin{tablenotes}
\small
\item \textit{Notes:} VAR(1) estimated by OLS on a balanced quarterly panel pooling four Middle East conflict episodes (1980Q3--1988Q3, 1990Q3--1991Q1, 2003Q1--2003Q2, 2019Q2--2019Q4); 188 observations. Standard errors and $t$-statistics from a residual bootstrap with $B = 500$ replications. Sovereign spreads from JPM EMBI+, Hilscher and Nosbusch (2010), \citet{borensztein2013debt}, and 5-year sovereign CDS post-2003. Standardized coefficients use the within-episode standardization; raw pass-through coefficients are reported on the original spread scale (basis points). See Online Appendix~A for full data construction.
\end{tablenotes}
\end{threeparttable}
\end{table}

\begin{figure}[htbp]
\centering
\includegraphics[width=0.85\textwidth]{figures/figure10_var_linkage.png}
\caption{Hand-Calibrated vs.\ VAR-Estimated Financial Linkage Matrix}
\label{fig:var}

\begin{minipage}{0.85\textwidth}
\vspace{0.3em}
\small \textit{Notes:} Heat-map comparison of the hand-calibrated linkage matrix used in the baseline (left panel) and the VAR(1) coefficients estimated on four historical Middle East conflict episodes (right panel). Diagonal entries are own-persistence; off-diagonal entries are spread-to-spread transmission coefficients. The two matrices are visually and quantitatively close.
\end{minipage}
\end{figure}

\subsection{Policy Counterfactuals}
\label{sec:counterfactuals}

To map out the global welfare frontier across alternative policy regimes, we run nine counterfactual experiments holding the war shock fixed at the baseline calibration. The experiments vary three margins: (i) China's stabilization intensity (baseline, no China, strong China); (ii) US monetary policy ($\phi_\pi \in \{1.0, 1.5, 2.5\}$, plus a coordinated easing scenario); and (iii) the financial channel (baseline, no financial frictions, no contagion, VAR-estimated linkage matrix). Welfare is computed in consumption-equivalent units following \citet{lucas1987models}.

Figure~\ref{fig:counterfactuals} reports ROW and China welfare losses across the nine regimes. Three results stand out and motivate Findings 1 and 2 in the Introduction. First, switching off financial frictions (E6) reduces global welfare loss by 4.88 percentage points---more than the entire trade-channel effect. Second, varying $\phi_\pi$ from 1.0 to 2.5 (E3 vs.\ E4) moves global welfare by only 0.11 percentage points, and even fully coordinated demand support (E5) cannot overcome the supply-side damage from the oil shock. Third, doubling the strength of China's stabilization (E8) reduces China's own welfare loss by 0.70 percentage points but reduces global welfare loss by only 0.18 percentage points, consistent with the binding-constraint interpretation in Section~\ref{sec:endogenous}. Together, these results sharpen the policy implications: the marginal returns to stabilization come from supply-side resilience (financial regulation, manufacturing capacity, energy security), not from demand management.

\begin{figure}[htbp]
\centering
\includegraphics[width=0.95\textwidth]{figures/figure11_counterfactuals.png}
\caption{Welfare Counterfactuals Across Nine Policy Regimes}
\label{fig:counterfactuals}

\begin{minipage}{0.95\textwidth}
\vspace{0.3em}
\small \textit{Notes:} Consumption-equivalent welfare loss (percent of permanent consumption, Lucas 1987) for ROW (left) and China (right) under nine policy experiments. E1: Baseline; E2: No China stabilization; E3: Aggressive Fed ($\phi_\pi = 2.5$); E4: Dovish Fed ($\phi_\pi = 1.0$); E5: Coordinated demand support; E6: No financial channels; E7: No financial contagion; E8: Strong China stabilization; E9: VAR-estimated linkage matrix. The welfare-loss reduction from E6 (4.88pp global) and the small spread between E3--E5 jointly motivate Findings 1 and 2.
\end{minipage}
\end{figure}


% ============================================================================
% VI. CONCLUSION
% ============================================================================

\section{Conclusion}
\label{sec:conclusion}

This paper develops a four-region New Keynesian DSGE model to analyze the macroeconomic consequences of a US--Israel conflict with Iran featuring a Hormuz Strait blockade. By embedding the trade-exposure spillover mechanism of \citet{federle2026price} within a structural framework augmented with financial channels, and by validating the structural results with a second-order solution, an empirically estimated linkage matrix, and a battery of policy counterfactuals, we deliver three contributions to the literature on the economics of war that bear directly on policy.

First, the Hormuz chokepoint creates a uniquely potent transmission mechanism that operates primarily through the financial channel. ROW output losses increase by a factor of 3.45$\times$ when financial contagion is incorporated, and the financial channel alone accounts for 4.88 percentage points of cumulative global welfare loss---more than the entire trade channel. A structural VAR(1) estimated on four historical Middle East conflict episodes confirms that Iran-to-third-country spread transmission is large ($+0.50$ to $+0.57$) and statistically significant ($t > 5$), validating the calibrated linkage matrix. The implication for macroprudential policy is concrete: financial regulation aimed at limiting cross-border spread spillovers would deliver welfare gains an order of magnitude larger than any plausible monetary-policy adjustment.

Second, conventional monetary policy is largely impotent in this rare-disaster setting. Across nine policy counterfactuals, varying Taylor-rule aggressiveness from $\phi_\pi = 1.0$ to $\phi_\pi = 2.5$ moves global welfare by less than 0.11 percentage points, and even fully coordinated demand support cannot offset the supply-side damage from a 63 percent oil price spike. The action is in fiscal coordination, energy security, and supply-chain redundancy, not in the central bank's reaction function. This finding sharpens the policy debate triggered by the recent literature on the limits of conventional monetary policy in disaster states \citep{barro2006rare}.

Third, China's manufacturing capacity, not its strategic petroleum reserves, is the binding stabilization channel. When deterministic stabilization paths are replaced by endogenous Taylor-rule-style reaction functions, the SPR rule binds at its physical capacity ceiling for the first eight quarters under any plausible parameterization, indicating that the marginal returns to additional reserves are zero. By contrast, doubling the manufacturing-reaction coefficient reduces China's own welfare loss by 0.70 percentage points. This finding challenges the prevailing narrative that frames China's global economic role primarily in terms of competitive disruption and reframes a familiar ``China shock'' literature \citep{autor2013china}: in the context of a major geopolitical crisis, China's diversified industrial base serves as a global supply-chain buffer with positive externalities for both China and its trading partners.

Methodologically, the second-order perturbation exercise quantifies---rather than dismisses---the bias from sequential linearization: Iran's peak loss deepens by 5.9 percentage points under the nonlinear solution, but the cross-regional spillover pattern and the qualitative ordering of all three findings are preserved. The VAR-estimated linkage matrix moves the headline welfare numbers by less than 5 percent. We interpret these robustness exercises as confirming that the structural conclusions are not artifacts of any single modeling choice.

The policy implications follow directly from the three findings. The Strait of Hormuz remains a critical vulnerability for the global economy, and military planners should incorporate the full general-equilibrium costs---dominated by financial contagion---into strategic assessments. Central banks should not expect conventional monetary policy to offset rare-disaster shocks; the policy lever lies elsewhere. And for China, the results suggest that investments in manufacturing resilience and supply-chain capacity generate positive global externalities that justify greater coordination than the existing literature acknowledges. Future work could embed our framework in an estimated Bayesian DSGE model with richer financial frictions, endogenous conflict dynamics, and heterogeneous agents, and could extend the empirical VAR by exploiting higher-frequency CDS data on additional historical episodes.


% ============================================================================
% REFERENCES
% ============================================================================

\newpage
\bibliographystyle{aer}

\begin{thebibliography}{99}

\bibitem[Abadie and Gardeazabal(2003)]{abadie2003economic}
Abadie, Alberto, and Javier Gardeazabal. 2003. ``The Economic Costs of Conflict: A Case Study of the Basque Country.'' \textit{American Economic Review}, 93(1): 113--132.

\bibitem[Autor, Dorn, and Hanson(2013)]{autor2013china}
Autor, David H., David Dorn, and Gordon H. Hanson. 2013. ``The China Syndrome: Local Labor Market Effects of Import Competition in the United States.'' \textit{American Economic Review}, 103(6): 2121--2168.

\bibitem[Barro(2006)]{barro2006rare}
Barro, Robert J. 2006. ``Rare Disasters and Asset Markets in the Twentieth Century.'' \textit{Quarterly Journal of Economics}, 121(3): 823--866.

\bibitem[Bernanke(2005)]{bernanke2005global}
Bernanke, Ben S. 2005. ``The Global Saving Glut and the U.S.\ Current Account Deficit.'' Sandridge Lecture, Virginia Association of Economists.

\bibitem[Blanchard and Gal\'{i}(2007)]{blanchard2007macroeconomic}
Blanchard, Olivier J., and Jordi Gal\'{i}. 2007. ``The Macroeconomic Effects of Oil Price Shocks: Why Are the 2000s So Different from the 1970s?'' NBER Working Paper 13368.

\bibitem[Bodenstein, Erceg, and Guerrieri(2011)]{bodenstein2011oil}
Bodenstein, Martin, Christopher J. Erceg, and Luca Guerrieri. 2011. ``Oil Shocks and External Adjustment.'' \textit{Journal of International Economics}, 83(2): 168--184.

\bibitem[Borensztein and Panizza(2013)]{borensztein2013debt}
Borensztein, Eduardo, and Ugo Panizza. 2013. ``The Costs of Sovereign Default.'' \textit{IMF Staff Papers}, 56(4): 683--741.

\bibitem[Caldara and Iacoviello(2022)]{caldara2022measuring}
Caldara, Dario, and Matteo Iacoviello. 2022. ``Measuring Geopolitical Risk.'' \textit{American Economic Review}, 112(4): 1194--1225.

\bibitem[{Energy Information Administration}(2023)]{eia2023hormuz}
Energy Information Administration. 2023. ``The Strait of Hormuz Is the World's Most Important Oil Transit Chokepoint.'' EIA Today in Energy, June 20.

\bibitem[Federle et~al.(2026)]{federle2026price}
Federle, Jonathan, Andr\'{e} Meier, Gernot J. M\"{u}ller, Willi Mutschler, and Moritz Schularick. 2026. ``The Price of War.'' \textit{American Economic Review}, 116(3): 791--827.

\bibitem[Gal\'{i}(2015)]{gali2015monetary}
Gal\'{i}, Jordi. 2015. \textit{Monetary Policy, Inflation, and the Business Cycle: An Introduction to the New Keynesian Framework}. 2nd ed. Princeton: Princeton University Press.

\bibitem[Hamilton(1983)]{hamilton1983oil}
Hamilton, James D. 1983. ``Oil and the Macroeconomy since World War II.'' \textit{Journal of Political Economy}, 91(2): 228--248.

\bibitem[Hamilton(2009)]{hamilton2009understanding}
Hamilton, James D. 2009. ``Understanding Crude Oil Prices.'' \textit{Energy Journal}, 30(2): 179--206.

\bibitem[Helveston et~al.(2019)]{helveston2019china}
Helveston, John P., Yimin He, and Michael R. Davidson. 2019. ``Quantifying the Cost Savings of Global Solar Photovoltaic Supply Chains.'' \textit{Nature}, 612: 83--87.

\bibitem[Hilscher and Nosbusch(2010)]{hilscher2010determinants}
Hilscher, Jens, and Yves Nosbusch. 2010. ``Determinants of Sovereign Risk: Macroeconomic Fundamentals and the Pricing of Sovereign Debt.'' \textit{Review of Finance}, 14(2): 235--262.

\bibitem[Jord\`{a}(2005)]{jorda2005estimation}
Jord\`{a}, \`{O}scar. 2005. ``Estimation and Inference of Impulse Responses by Local Projections.'' \textit{American Economic Review}, 95(1): 161--182.

\bibitem[Kilian(2009)]{kilian2009not}
Kilian, Lutz. 2009. ``Not All Oil Price Shocks Are Alike: Disentangling Demand and Supply Shocks in the Crude Oil Market.'' \textit{American Economic Review}, 99(3): 1053--1069.

\bibitem[Longstaff et~al.(2011)]{longstaff2011financial}
Longstaff, Francis A., Jun Pan, Lasse H. Pedersen, and Kenneth J. Singleton. 2011. ``How Sovereign Is Sovereign Credit Risk?'' \textit{American Economic Journal: Macroeconomics}, 3(2): 75--103.

\bibitem[Lucas(1987)]{lucas1987models}
Lucas, Robert E., Jr. 1987. \textit{Models of Business Cycles}. Oxford: Basil Blackwell.

\bibitem[Ramey(2011)]{ramey2011identifying}
Ramey, Valerie A. 2011. ``Identifying Government Spending Shocks: It's All in the Timing.'' \textit{Quarterly Journal of Economics}, 126(1): 1--50.

\bibitem[Rey(2015)]{rey2015dilemma}
Rey, H\'{e}l\`{e}ne. 2015. ``Dilemma Not Trilemma: The Global Financial Cycle and Monetary Policy Independence.'' NBER Working Paper 21162.

\bibitem[Zhang et~al.(2020)]{zhang2020china}
Zhang, Dayong, Qiang Ji, and Apostolos Kiohos. 2020. ``China's Oil Strategic Petroleum Reserve: A DSGE Analysis.'' \textit{Energy Economics}, 88: 104748.

\end{thebibliography}


% ============================================================================
% APPENDIX (FOR ONLINE PUBLICATION)
% ============================================================================

\newpage
\appendix
\section*{Appendix: Model Equations Summary}
\addcontentsline{toc}{section}{Appendix}

For reference, we collect the complete system of equations for each region $i$:

\begin{align}
\pi_{i,t} &= \beta \, \mathbb{E}_t[\pi_{i,t+1}] + \kappa_i \, \hat{y}_{i,t-1} + \xi^{\text{supply}}_{i,t} \tag{A.1} \\[6pt]
\hat{y}_{i,t} &= \rho \, \hat{y}_{i,t-1} - \frac{1}{\sigma_i}(i_{i,t} - \mathbb{E}_t[\pi_{i,t+1}]) + 0.8 \, g^{\text{mil}}_{i,t} - \tau^{\text{trade}}_{i,t} - 0.3 \, \alpha_{e,i} \Delta p^{\text{oil}}_t + \Phi^{\text{fin}}_{i,t} \tag{A.2} \\[6pt]
i_{i,t} &= \rho_{i,i} \, i_{i,t-1} + (1-\rho_{i,i})(\phi_{\pi,i} \pi_{i,t} + \phi_{y,i} \hat{y}_{i,t-1}) \tag{A.3} \\[6pt]
s_{i,t} &= s^{\text{shock}}_{i,t} + 0.3 \max(-\hat{y}_{i,t-1}, 0) + 0.2 \, s_{i,t-1} \tag{A.4} \\[6pt]
k_{i,t} &= k^{\text{shock}}_{i,t} - 2 s_{i,t} + 0.5(i_{i,t} - 0.02) + 0.3 \, k_{i,t-1} \tag{A.5} \\[6pt]
q_{i,t} &= 0.6 \, q_{i,t-1} + 10 g^e_{i,t} - 8(i_{i,t} + s_{i,t}) - 3 \max(-\hat{y}_{i,t},0) + q^{\text{shock}}_{i,t} \tag{A.6} \\[6pt]
e_{i,t} &= 0.4 \, e_{i,t-1} + 0.5 \, s_{i,t} - 0.3 \, k_{i,t} + \iota^{\text{fx}}_{i,t} \tag{A.7} \\[6pt]
\Phi^{\text{fin}}_{i,t} &= -0.15 \, s_{i,t} - 0.10 \max(-k_{i,t}, 0) \tag{A.8}
\end{align}

The global oil market clears according to:
\begin{align}
\Delta S_t &= \lambda^{\text{block}}_t (s^H - s^{\text{Iran}}) + s^{\text{Iran}} \lambda^{\text{prod}} - 0.3 \lambda^{\text{block}}_t (s^H - s^{\text{Iran}}) \tag{A.9} \\[6pt]
\frac{\Delta P^{\text{oil}}_t}{P^{\text{oil}}} &= \frac{\Delta S_t - \text{SPR}_t + \Delta D_t}{|\varepsilon_d| + \varepsilon_s} \tag{A.10}
\end{align}

Financial contagion operates through:
\begin{equation}
\tilde{s}^{\text{shock}}_{j,t} = s^{\text{shock}}_{j,t} + 0.1 \Big(\textstyle\sum_{i \neq j} L_{ij} s_{i,t-1}\Big)(1 + 2\max(-k_{\text{Iran},t}, 0)) \tag{A.11}
\end{equation}

\end{document}
